\documentclass[11pt,a4paper]{article}
\usepackage{fontspec}
\usepackage{polyglossia}
\setdefaultlanguage{russian}
\setotherlanguage{english}
\setmainfont{Times New Roman}
\setsansfont{Arial}
\setmonofont{Consolas}
\usepackage[margin=2.2cm]{geometry}
\usepackage{amsmath,amssymb,amsthm,mathtools}
\usepackage{booktabs}
\usepackage{enumitem}
\usepackage[colorlinks=true,linkcolor=blue!60!black,citecolor=blue!60!black,urlcolor=blue!60!black]{hyperref}

\newtheorem{theorem}{Теорема}[section]
\newtheorem{corollary}[theorem]{Следствие}
\newtheorem{proposition}[theorem]{Предложение}
\theoremstyle{definition}
\newtheorem{definition}[theorem]{Определение}
\theoremstyle{remark}
\newtheorem{remark}[theorem]{Замечание}

\title{\textbf{Прикладная инженерия дискретной квантовой механики (ДКМ):\\ теоретическая электротехника фазовой эстафеты ($V=\mathcal{C}_3,\ I=\mathcal{S}_3$), беспульсационные квадратурные выпрямители, физическая реализация вычислений без тепла Ландауэра и кремниевая архитектура (\texttt{Ring9} / \texttt{Bus36} / \texttt{VerletCell72}) в Lean~4}}
\author{Д.\,А.~Каратаев\thanks{AI-Guiders Research (\url{https://github.com/AI-Guiders}). Формальная верификация в Lean~4 и аналитическая разработка выполнены при участии ИИ-партнёра \emph{Луч} (\emph{Luch}, AI-Guiders).}}
\date{2026}

\begin{document}
\maketitle

\begin{abstract}
Когда физика, химия, биология и физическая реализация вычислений объединяются в единое Естествознание на основе дискретной квантовой механики (ДКМ) и двусортной структуры носителей $\langle\mathcal{M},\mathcal{T}\rangle$ ($\mathcal{M} = \mathbb{Z}_{512}$ на оболочке кубооктаэдра $O_2(3)$ при $1+2^3=3^2$), электротехника и вычислительная техника сводятся к точному целочисленному согласованию фазовых импедансов. Во-первых, доказано, что в гранецентрированной кубической (ГЦК) решётке металлических проводников ($\mathrm{Cu, Ag, Au, Al}$), чья первая координационная сфера из $12$ ближайших соседей совпадает с кубооктаэдром $O_2(3)$, электрический ток представляет собой торсионную эстафету вращения фазы $I(\Delta\varphi) = \mathcal{S}_3(\Delta\varphi)$, сопряжённую с потенциальным натяжением $V(\Delta\varphi) = \mathcal{C}_3(\Delta\varphi) = 128 - |\Delta\varphi|$. Из инварианта грани $|\mathcal{C}_3(\Delta\varphi)| + |\mathcal{S}_3(\Delta\varphi)| \equiv 128$ выводятся пять отдельных теорем (\texttt{thm5} в \texttt{LogicCohomology.lean}): нагрузочная прямая постоянного тока ($P(\Delta\varphi) \le 64^2 = 4096$), законы Кирхгофа и Фарадея ($d_1\circ d_0 \equiv 0$, $d_2\circ d_1 \equiv 0$), ортогональность реактивной мощности и противофазное гашение в $LC$-резонансе ($\mathcal{C}_3(\varphi+128) + \mathcal{C}_3(\varphi+384) \equiv 0$), однополупериодный, мостовой выпрямитель Гретца и бесконденсаторный квадратурный выпрямитель с $0\%$ пульсаций ($|\mathcal{C}_3(\varphi)| + |\mathcal{S}_3(\varphi)| \equiv 128$), а также трансформаторное преобразование импеданса $Z_{\mathrm{in}} = (N_1/N_2)^2 Z_{\mathrm{load}}$. Во-вторых, формализована физическая реализация информатики в кремнии и акусто-оптической аппаратуре: обратимые вычисления в коридоре Гейзенберга $|\Delta\varphi|\le 40 < 41$ с нулевым выделением тепла Ландауэра ($\hat{n}_E = 0$), диагонализация взаимной индукции $M_{jk} = 0$ между $12$ катушками $O_2(3)$ и безветвевая архитектура (\texttt{Ring9}, \texttt{Bus36}, \texttt{VerletCell72}).
\end{abstract}

\section{Теоретическая электротехника фазовой эстафеты на ГЦК-кубооктаэдрах (\texttt{LogicCohomology.lean})}
\label{sec:electrical-engineering}

В классической теории Друде скорость дрейфа электронов в медном проводе составляет $v_{\mathrm{drift}} \sim 10^{-4}\;\text{м/с}$, что вступает в противоречие с передачей электрической мощности со скоростью света в линии. В ДКМ на дискретном носителе $\mathcal{M} = \mathbb{Z}_{512}$ (в теоретико-модельном смысле А.\,И.~Мальцева~\cite{Malcev1970}) кристаллическая решётка ГЦК-проводников ($\mathrm{Cu},\ Z=29$; $\mathrm{Ag},\ Z=47$; $\mathrm{Au},\ Z=79$; $\mathrm{Al},\ Z=13$) образована атомными узлами, чья первая координационная сфера из $12$ ближайших соседей совпадает с кубооктаэдром $O_2(3)$.

\begin{theorem}[Нагрузочная прямая постоянного тока, согласование импедансов (MPPT) и сверхпроводимость против джоулева нагрева, \texttt{thm5}]
\label{thm:dc-ohm-mppt}
Стационарный сдвиг фазы $\Delta\varphi = k \in [0,128]$ между соседними узлами кубооктаэдра делит бюджет четверти кольца $Q_\square(3) = 128$ между потенциальным напряжением $V(k) = \mathcal{C}_3(k) = 128 - k$ и током фазовой эстафеты $I(k) = \mathcal{S}_3(k) = k$:
\begin{equation}
  \label{eq:dc-ohm-mppt}
  V(k) + I(k) \;\equiv\; 128, \qquad
  P(k) \;=\; V(k)\,I(k) \;=\; 64^2 - (64 - k)^2 \;\le\; 4\,096 \;=\; K_P(3)^2,
\end{equation}
достигая максимума передачи мощности при $k^* = K_P(3) = 64$ ($V = I = 64$). Внутри подпорогового коридора $|k| \le 40 < 41$ излучение фононов тождественно равно нулю ($\hat{n}_E(k) = \lfloor|k|/41\rfloor = 0$, баллистический/сверхпроводящий режим); при $|k| \ge 41$ на неоднородностях решётки излучается $\hat{n}_E(k) \in \{1,\dots,6\}$ фононных квантов (закон Джоуля--Ленца) с рождением вихревых магнитных плакет $B_f = 512\,n_f$.
\end{theorem}

\begin{theorem}[Законы Кирхгофа (KVL и KCL) и индукции Фарадея из нильпотентности кограницы де~Рама, \texttt{thm5}]
\label{thm:kirchhoff-faraday}
Вокруг любого замкнутого контура $\partial f$ тождество $(d_1 d_0\varphi)_f \equiv 0$ даёт второй закон Кирхгофа в потенциальных цепях и закон индукции Фарадея $(d_1 r)_f = -512(d_1 q)_f = 512\,n_f$ при сцеплении с вихревым потоком, тогда как $d_2\circ d_1 \equiv 0$ и $\sum_{k=0}^{11}\vec{e}_k = \mathbf{0}$ дают первый закон Кирхгофа в каждом узле.
\end{theorem}

\begin{theorem}[Ортогональность реактивной мощности переменного тока, энергия Парсеваля и резонанс $LC$-контура, \texttt{thm5}]
\label{thm:ac-rlc-laws}
Ёмкостный ($+128$ тиков) и индуктивный ($-128 \equiv +384$ тика) сдвиги фазы на $\mathbb{Z}_{512}$ удовлетворяют
\begin{equation}
  \label{eq:ac-rlc-laws}
  \sum_{\varphi=0}^{511} \mathcal{C}_3(\varphi)\,\mathcal{S}_3(\varphi) = 0, \qquad
  \sum_{\varphi=0}^{511} \mathcal{C}_3(\varphi)^2 = 2\,796\,288, \qquad
  \mathcal{C}_3(\varphi + 128) + \mathcal{C}_3(\varphi + 384) \;\equiv\; 0 \quad (\forall\varphi\in\mathbb{Z}_{512}),
\end{equation}
доказывая равенство нулю работы реактивной мощности за период и поточечную взаимную аннигиляцию реактивных напряжений $L$ и $C$ в резонансе.
\end{theorem}

\begin{theorem}[Однополупериодный диод, двухполупериодный мост Гретца и квадратурный выпрямитель с $0\%$ пульсаций, \texttt{thm5}]
\label{thm:rectifiers}
Безветвевой однополупериодный выпрямитель $\mathrm{rect}_{1/2}(x) = (|x|+x)/2$ и двухполупериодный $4$-диодный мост Гретца $|x|$ дают точные целочисленные средние токи на $\mathbb{Z}_{512}$ без иррациональных множителей $\pi$:
\begin{equation}
  \label{eq:rectifiers}
  \sum_{\varphi=0}^{511}\mathrm{rect}_{1/2}\bigl(\mathcal{S}_3(\varphi)\bigr) = 16\,384 = 128^2 \;\;(\bar{I}_{1/2} = 32), \qquad
  \sum_{\varphi=0}^{511}\bigl|\mathcal{S}_3(\varphi)\bigr| = 32\,768 = 2\times 128^2 \;\;(\bar{I}_{\mathrm{Graetz}} = 64 = K_P),
\end{equation}
а двухфазный квадратурный мостовой выпрямитель ($\varphi$ и $\varphi - 128$) обеспечивает постоянный выход без конденсатора с $0\%$ пульсаций на каждом тике: $|\mathcal{C}_3(\varphi)| + |\mathcal{S}_3(\varphi)| \equiv 128$.
\end{theorem}

\begin{theorem}[Идеальный трансформатор и квадратичное преобразование импеданса]
\label{thm:transformer}
Для обмоток $N_1 : N_2$ на общем замкнутом магнитопроводе ($\Phi_{\mathrm{core}} = 512\,n_f$) имеем $V_1 = N_1\varepsilon_1$, $V_2 = N_2\varepsilon_1$, $I_1 = N_2 i_0$, $I_2 = N_1 i_0$, откуда $N_2 V_1 - N_1 V_2 \equiv 0$, $N_1 I_1 - N_2 I_2 \equiv 0$, $V_1 I_1 \equiv V_2 I_2$ и $Z_{\mathrm{in}} = (N_1/N_2)^2 Z_{\mathrm{load}}$.
\end{theorem}

\section{Физическая реализация информатики: кремниевая и акусто-оптическая архитектура без тепла Ландауэра}
\label{sec:silicon-hardware}

В то время как абстрактная информатика (полнота Поста и фазовая редукция 3-SAT, \texttt{thm4} в \texttt{LogicCohomology.lean}) принадлежит чистой математике, её физическая реализация во Вселенной без внешней ленты Тьюринга определяется инженерными законами ДКМ:
\begin{enumerate}[label=\textbf{\arabic*.},leftmargin=*]
\item \textbf{Подпороговые обратимые вычисления без тепла Ландауэра ($|\Delta\varphi|\le 40 < 41$):}
  поскольку двухтактовый симплектический шаг Верле $\varphi(t+1) = (2\varphi(t) - \varphi(t-1) + \Phi_{\mathrm{gate}})\bmod 512$ биективен на $72$-битном регистре ячейки (\texttt{VerletCell72}), любая логическая передача внутри мёртвой зоны Гейзенберга $|\Phi_{\mathrm{gate}}|\le 40$ излучает $\hat{n}_E = 0$ тепловых квантов, обходя предел стирания Ландауэра $k_B T\ln 2$.
\item \textbf{Диагонализация взаимной индукции между 12 катушками $O_2(3)$:}
  вокруг любой выбранной оси катушки $\vec{e}_0\in O_2(3)$ ($\|\vec{e}_0\|_2^2 = 2$) остальные $11$ катушек распадаются на $4$ соседа под углом $60^\circ$ ($\vec{e}_0\cdot\vec{e}_j = +1$), $2$ ортогональных соседа под углом $90^\circ$ ($\vec{e}_0\cdot\vec{e}_j = 0$), $4$ тупоугольных соседа под углом $120^\circ$ ($\vec{e}_0\cdot\vec{e}_j = -1$) и $1$ антиподальный партнёр под углом $180^\circ$ ($\vec{e}_0\cdot(-\vec{e}_0) = -2$). Внеосевая взаимная индукция сокращается тождественно: $4(+1) + 2(0) + 4(-1) = 0$.
\item \textbf{Безветвевой аппаратный стек на кремнии (\texttt{Verilog} и \texttt{.NET~11}: \texttt{Ring9}, \texttt{Bus36}, \texttt{VerletCell72}):}
  реализован в \texttt{hardware/cuboctahedron\_carrier\_asic.v} (18 синтезируемых модулей Verilog) и \texttt{CarrierCore.cs} (\url{https://github.com/AI-Guiders/guiders-physics}):
  \begin{itemize}[noitemsep]
  \item \texttt{Ring9} ($9$-битная фазовая арифметика с бабочкой ротатора Грея $\mathcal{C}_3,\mathcal{S}_3$ и делением Безу);
  \item \texttt{Bus36} ($4\times 9 = 36$-битная четырёхполосная шина по основанию $512$ без погрешностей плавающей запятой);
  \item \texttt{Module~17 \& 18} (безветвевое восстановление всех $512/512$ фаз из акустического давления $P_{\mathrm{sound}} = \mu_0\mu_1(128-\mathcal{C}_3)$ и хирального тока $J_{\mathrm{rxn}} = \mu_0\mu_1\mathcal{S}_3$, 1-тактовая замкнутая ФАПЧ и 4-канальный нильпотентный томограф Неймана $(2I_4 + N_{\mathrm{nil}})^{-1} = \frac{1}{24}(12 I_4 - 6 N_{\mathrm{nil}} + 3 N_{\mathrm{nil}}^2)$ при $N_{\mathrm{nil}}^3 \equiv 0$).
  \end{itemize}
\end{enumerate}

\section{Единый ньютоновский фазовый гомеоморфизм всех инженерных дисциплин (\texttt{HardwarePhotovoltaics.lean})}
\label{sec:unified-engineering-homeomorphism}

Поскольку закон эволюции дискретного носителя $\mathcal{M} = O_2(3)\times\mathbb{Z}_{512}$ есть второй закон Ньютона в точных целых числах ($\Delta^2\varphi(t) = \varphi(t+1)-2\varphi(t)+\varphi(t-1) = \lfloor N_{\mathrm{torque}}/D(\mathbf{x})\rfloor$) в сочетании с первым законом Ньютона в зазоре зацепления ($|\Delta\varphi|\le 40 < 41 \implies \hat{n}_E = 0$) и третьим законом Ньютона вдоль всех $12$ связей кубооктаэдра ($\mathcal{S}_3(-\Delta\varphi) = -\mathcal{S}_3(\Delta\varphi)$), \textbf{все шесть классических инженерных дисциплин сворачиваются в единую систему фазового согласования импедансов, управляемую пятью каноническими порогами кольца} ($|\Delta\varphi|\le 40 < 41$, $K_P = 64$, $Q_\square = 128$, $\Phi_{\mathrm{Pauli}} = 256$), как показано в Таблице~\ref{tab:engineering-homeomorphism}.

\begin{table}[htbp]
\centering
\caption{Сведение шести инженерных дисциплин от континуальных аппроксимаций $\mathcal{T}$-слоя к целочисленной ньютоновской фазовой механике на $O_2(3)\times\mathbb{Z}_{512}$.}
\label{tab:engineering-homeomorphism}
\footnotesize
\begin{tabular}{@{}p{0.21\textwidth}p{0.35\textwidth}p{0.39\textwidth}@{}}
\toprule
\textbf{Отрасль инженерии} & \textbf{Ограничение континуального $\mathcal{T}$-слоя ($\mathbb{R}^3$)} & \textbf{Решение в ньютоновской $\mathcal{M}$-механике ($O_2(3)\times\mathbb{Z}_{512}$)} \\
\midrule
\textbf{1. Электротехника и силовая электроника} &
Гармоники, циркуляция реактивной мощности, пульсации конденсаторов и джоулевы потери $I^2 R$. &
$V=\mathcal{C}_3,\ I=\mathcal{S}_3$ на $L^1$-ромбе: $|V|+|I|\equiv 128$ ($0\%$ пульсаций без конденсатора), MPPT при $k^*=64$ ($P_{\max}=4\,096$), гашение взаимной индукции $12$ катушек $4(+1)+2(0)+4(-1)=0$ и $\hat{n}_E=0$ при $|\Delta\varphi|\le 40$. \\
\addlinespace[3pt]
\textbf{2. Полупроводники и фотовольтаика} &
Предел Шокли--Квейссера $\sim 33\%$ ($60\%$ потерь синих фотонов в тепло) и оже-рекомбинация. &
Синглетное деление $441\;\text{нм}\to 2\times 882\;\text{нм} < 1\,107\;\text{нм}$ при шаге $21 < 41$ ($\hat{n}_E=0$); кислородный замок седловины $Z(\mathrm{C_{18}H_{12}})+Z(\mathrm{O})=120+8=128=Q_\square$ и барьер Паули $\mathrm{SiO}_4$ $46^2=2\,116>2\,048$. \\
\addlinespace[3pt]
\textbf{3. Вычислительная техника и защита} &
Тепло необратимых булевых вентилей $\ge k_{\mathrm{B}}T\ln 2$ и зависимость криптографии от гипотезы $P\stackrel{?}{=}NP$. &
Обратимое АЛУ Верле Канала~1 (\texttt{Ring9/Bus36/VerletCell72}) с четверть-квадратным умножением $(x+y)^2-(x-y)^2=4xy$ при $\hat{n}_E=0$; барьеры решателя внутри носителя ($72^2=5\,184$, $7^T$) и встроенный шифр Вернама $H(M\mid T)=H(M)$. \\
\addlinespace[3pt]
\textbf{4. Теплотехника и термодинамика} &
Трактует тепло как необратимый хаос, требующий внешнего холодильника или стирания памяти. &
$12$-лучевой расщепитель кубооктаэдра без памяти делит даже $6$-квантовый удар Паули $\Phi_{\mathrm{Pauli}}=256$ на $\lfloor 256/12\rfloor = 21 < 41$, переводя тепловой удар в баллистический ток ($12\times 21=252>0$) при $\hat{n}_E(21)=0$. \\
\addlinespace[3pt]
\textbf{5. Гидрогазодинамика и акустика} &
Уравнения Навье--Стокса дают континуальные сингулярности и неквантованную турбулентность. &
$4$-канальный нильпотентный каскад $\mathcal{M}_{\mathrm{chan}} = I_4 + N_{\mathrm{nil}}$ ($N_{\mathrm{nil}}^3=0$, $\det(2\mathcal{M})=24$) точно обращается полиномом $I_4 - N_{\mathrm{nil}} + N_{\mathrm{nil}}^2$; турбулентность квантуется вихрями $512 n_f$ выше порога $41$, критическое число Маха $c_0/c=\sqrt{2}$. \\
\addlinespace[3pt]
\textbf{6. Сопромат, катализ и ТММ приводов} &
Эмпирические коэффициенты усталости, накопление ошибки Рунге--Кутты и фазовый джиттер каскадов. &
Упругий режим без усталости при $|\Delta\varphi|\le 40 < 41$; седловина катализа в верхней мёртвой точке $Q_\square=128$ (коллапс барьера в $587$ раз); потолок Паули $D_{\max}=833\varepsilon$ при $256$; точный привод Верле $\Delta^2\varphi\equiv\Phi_{\mathrm{imp}}$ и гашение джиттера $(\varphi_{\mathrm{target}}-\Sigma j)+\Sigma j = \varphi_{\mathrm{target}}$. \\
\bottomrule
\end{tabular}
\end{table}

\begin{thebibliography}{9}
\bibitem{Malcev1970}
А.\,И.~Мальцев, \emph{Алгебраические системы}, М.: Наука, 1970 (англ.\ пер.: Springer-Verlag, 1973).
\bibitem{KarataevRotator1}
Д.\,А.~Каратаев, «Ротаторное исчисление на кубооктаэдре: дискретная экспонента Эйлера на $\mathbb{Z}_{512}$, универсальный константный вронскиан и замкнутые квадратуры на кубооктаэдре $O_2(3)$ в Lean~4», \url{https://github.com/AI-Guiders/scientific-papers}, 2026.
\bibitem{Moura2021}
L.~de~Moura and S.~Ullrich, ``The Lean~4 theorem prover and programming language,'' in \emph{CADE~28}, LNCS vol.~12699, Springer, 2021, pp.~625--635.
\end{thebibliography}

\end{document}
